
% alternative formatting for algorithm
\iffalse
\begin{algo}
The procedure
\(\texttt{Access}(\texttt{op},\texttt{addr},\texttt{data}^{*})\)
where \(\texttt{op} = \texttt{read}\) or
\(\texttt{op} = \texttt{write}\)

\textbf{Assume:} Each block is of the form
\(( \texttt{addr},\texttt{data},l )\) where
\(l\) denotes the block's current designated path.

\begin{enumerate}
\item
  Let \(l^{*}\) be a random value from \(1\) to \(N\). Assign
  \(l \leftarrow\texttt{position}\left\lbrack \texttt{addr} \right\rbrack\),
  \(\texttt{position}\left\lbrack \texttt{addr} \right\rbrack \leftarrow l^{*}\).
\item
  For each bucket from leaf \(l\) to root do:

  \begin{enumerate}
  \setcounter{enumii}{2}
  \item
    Scan \(\texttt{bucket}\), and if
    \(( \texttt{addr},\texttt{data}_{0},\_ ) \in \texttt{bucket}\)
    then let
    \(\texttt{data}^{*} \leftarrow \texttt{data}_{0}\) and
    remove this block from \(\texttt{bucket}\).
  \end{enumerate}
\item
  End for
\item
  If \(\texttt{op} = \texttt{read}\) then add
  \(( \texttt{addr},\texttt{data}^{*},l^{*} )\)
  to the root bucket; else add
  \(( \texttt{addr},\texttt{data},l^{*} )\) to
  the root bucket.
\item
  Call the \(\texttt{evict}\) subroutine.
\item
  Return \(\texttt{data}^{*}\).
\end{enumerate}
\end{algo}

\begin{algo}
\phantomsection\label{alg:evict}
strong(Algorithm):\\
The procedure \(\texttt{evict}\)

\begin{enumerate}
\item
  For each level \(d\) from the root to the level of leaves \(- 1\), do:

  \begin{enumerate}
  \setcounter{enumii}{1}
  \item
    \(\texttt{bucket}_{0}\),
    \(\texttt{bucket}_{1} \leftarrow\) randomly choose \(2\)
    distinct buckets in the level \(d\) (for the root level, pick one
    bucket).
  \item
    For
    \(\texttt{bucket} \in \{\texttt{bucket}_{0},\texttt{bucket}_{1}\}\)
    do:

    \begin{enumerate}
    \setcounter{enumiii}{3}
    \item
      \(\mathsf{\text{block}} :=\) pop a real block from
      \(\mathsf{\text{bucket}}\) if one exists; else
      \(\mathsf{\text{block}} := (\bot,\bot,\bot)\).
    \item
      For each of the two children of \(\texttt{bucket}\) in a
      fixed order: scan the child bucket reading and writing every
      block. If \(\mathsf{\text{block}}\) is real and wants to go to the
      child, write \(\mathsf{\text{block}}\) to an empty slot in the
      child bucket.
    \end{enumerate}
  \item
    End for
  \end{enumerate}
\item
  End for
\end{enumerate}
\end{algo}
\fi
